% ---------------------------------------------------------------------------
% CVPR / ICCV Conference Paper Template
%
% Usage:
%   1. Download official cvpr.sty from: https://cvpr.thecvf.com/author_guidelines
%   2. Place cvpr.sty in this directory
%   3. Compile: pdflatex template.tex && bibtex template && pdflatex template.tex x2
%
% Conference: CVPR (Computer Vision and Pattern Recognition) / ICCV (International Conference on Computer Vision)
% Page limit: 8 pages (CVPR) + references
% Typical structure follows CVPR 2-column format
% ---------------------------------------------------------------------------

\documentclass[10pt,twocolumn,letterpaper]{article}

% --- UNCOMMENT when cvpr.sty is available ---
% \usepackage{cvpr}
% \usepackage[pagebackref,breaklinks,colorlinks,citecolor=blue]{hyperref}
% \cvprfinalcopy

\usepackage[utf8]{inputenc}
\usepackage[T1]{fontenc}
\usepackage{graphicx}
\usepackage{amsmath,amssymb}
\usepackage{booktabs}
\usepackage{multirow}
\usepackage{xcolor}
\usepackage[pagebackref,breaklinks,colorlinks,citecolor=blue]{hyperref}
\usepackage[margin=1in]{geometry}  % Remove when using cvpr.sty

% --- Paper Information ---
% \def\cvprPaperID{****}
% \def\confName{CVPR}
% \def\confYear{2026}

\begin{document}

% ===========================================================================
% TITLE
% ===========================================================================
\title{Your Paper Title Here}

\author{
  Author One${}^1$,
  Author Two${}^1$,
  Author Three${}^2$ \\
  \\
  ${}^1$Institution One, City, Country \\
  ${}^2$Institution Two, City, Country \\
  \texttt{\{email1, email2\}@institution.edu}
}

\maketitle

% ===========================================================================
% ABSTRACT
% ===========================================================================
\begin{abstract}
  Your abstract here. Follow the 2-6-2 structure:
  2 sentences background/gap,
  6 sentences method (First/Second/Finally),
  2 sentences results/implication.
  Target: $\sim$210 English words.

  \textbf{Example:} Establishing fine-grained semantic correspondence between
  natural language and visual content is the central challenge in cross-modal
  retrieval. Existing methods [gap]... First, we [component 1 + why].
  Second, we [component 2 + why]. Finally, we [component 3 + why].
  On [dataset], [Method] achieves [metric] = [value]\%, outperforming
  [baseline] by [X] percentage points.
\end{abstract}

% ===========================================================================
% 1. INTRODUCTION
% ===========================================================================
\section{Introduction}
\label{sec:intro}

% Para 1: Task importance and real-world applications
% Para 2: Current progress -- 2-3 method families with specific examples
% Para 3: Remaining gap -- what existing methods fail to do and WHY
% Para 4: Proposed method teaser -- 2-3 sentences, WHAT + intuition
% Para 5: Contributions -- 3-4 bullets with evidence pointers

Your introduction text here.

Our main contributions are:
\begin{enumerate}
  \item Contribution 1 with evidence pointer (Table/Figure/Section).
  \item Contribution 2 with evidence pointer.
  \item Contribution 3 with evidence pointer.
\end{enumerate}

% ===========================================================================
% 2. RELATED WORK
% ===========================================================================
\section{Related Work}
\label{sec:related}

% Organize by themes, NOT paper-by-paper:
% 2.1 Theme A (2-3 papers, comparison to your method)
% 2.2 Theme B (2-3 papers, comparison to your method)
% 2.3 Theme C (gap evidence)

\subsection{Theme A: [Method Family Name]}

Discussion of related work organized by themes.

\subsection{Theme B: [Another Method Family]}

Discussion of related work.

% ===========================================================================
% 3. PROPOSED METHOD
% ===========================================================================
\section{Proposed Method}
\label{sec:method}

% 3.1 Overview -- architecture diagram reference (Fig.~1)
\subsection{Overview}
\label{sec:method-overview}

Figure~\ref{fig:architecture} illustrates the overall architecture.

\begin{figure}[t]
  \centering
  \fbox{\parbox{0.8\linewidth}{\centering
  \vspace{2cm}
  {\large [Architecture Diagram -- Fig.~1]}\\
  Insert architecture diagram here.
  \vspace{2cm}}}
  \caption{Overview of the proposed method.}
  \label{fig:architecture}
\end{figure}

% 3.2 Module B -- Input → Process → Output → Rationale
\subsection{Module B: [Component Name]}
\label{sec:method-moduleb}

Describe your first key component. For each module, follow this structure:
\begin{enumerate}
  \item \textbf{Input:} What goes in (tensor shapes, data types)
  \item \textbf{Process:} What computation happens (equations, algorithms)
  \item \textbf{Output:} What comes out
  \item \textbf{Rationale:} Why this design choice
\end{enumerate}

% 3.3 Module C
\subsection{Module C: [Component Name]}
\label{sec:method-modulec}

Describe your second key component.

% 3.4 Loss Function / Training Objective
\subsection{Training Objective}
\label{sec:method-loss}

The overall training objective combines:
\begin{equation}
  \mathcal{L}_{\text{total}} = \mathcal{L}_{\text{primary}} + \lambda \mathcal{L}_{\text{auxiliary}}
  \label{eq:loss}
\end{equation}

% ===========================================================================
% 4. EXPERIMENTS
% ===========================================================================
\section{Experiments}
\label{sec:experiments}

% 4.1 Datasets and Implementation Details
\subsection{Datasets and Implementation Details}
\label{sec:exp-setup}

\textbf{Datasets.} Describe datasets: name, size, splits, preprocessing.

\textbf{Implementation Details.} Include: framework, optimizer, learning rate,
scheduler, batch size, epochs, hardware, seeds, data augmentation.

\textbf{Baselines.} List all baselines with citations.

\textbf{Evaluation Metrics.} Define metrics with direction (higher/lower is better).

% 4.2 Comparison with State-of-the-Art
\subsection{Comparison with State-of-the-Art}
\label{sec:exp-main}

\begin{table}[t]
  \centering
  \caption{Main results comparison on [Dataset]. Best in \textbf{bold}, second-best \underline{underlined}.}
  \label{tab:main-results}
  \begin{tabular}{lccc}
    \toprule
    \textbf{Method} & \textbf{Metric 1} & \textbf{Metric 2} & \textbf{Metric 3} \\
    \midrule
    Baseline 1 \cite{ref1} & 00.0 & 00.0 & 00.0 \\
    Baseline 2 \cite{ref2} & 00.0 & 00.0 & 00.0 \\
    Baseline 3 \cite{ref3} & 00.0 & 00.0 & 00.0 \\
    \textbf{Ours}          & \textbf{00.0} & \textbf{00.0} & \textbf{00.0} \\
    \bottomrule
  \end{tabular}
\end{table}

% 4.3 Ablation Study
\subsection{Ablation Study}
\label{sec:exp-ablation}

\begin{table}[t]
  \centering
  \caption{Ablation study on [Dataset]. Each row removes one component from the full model.}
  \label{tab:ablation}
  \begin{tabular}{lcc}
    \toprule
    \textbf{Variant} & \textbf{Metric 1} & \textbf{$\Delta$} \\
    \midrule
    Full Model & 00.0 & -- \\
    w/o Component B & 00.0 & -0.0 \\
    w/o Component C & 00.0 & -0.0 \\
    w/o Both & 00.0 & -0.0 \\
    \bottomrule
  \end{tabular}
\end{table}

% 4.4 Efficiency Analysis (if applicable)
\subsection{Efficiency Analysis}
\label{sec:exp-efficiency}

% 4.5 Qualitative Analysis
\subsection{Qualitative Analysis}
\label{sec:exp-qualitative}

% ===========================================================================
% 5. CONCLUSION
% ===========================================================================
\section{Conclusion}
\label{sec:conclusion}

% Restate problem, method, key findings.
% Do NOT add new claims or citations.
% End with 1-2 specific future work directions.

In this paper, we have presented [Method], which [problem + core mechanism].
Experiments on [N] datasets demonstrate that [key finding 1] and [key finding 2].
Limitations include [specific limitation]. Future work may explore [specific direction].

% ===========================================================================
% ACKNOWLEDGMENTS (optional)
% ===========================================================================
% \section*{Acknowledgments}
% Funding and thanks here.

% ===========================================================================
% REFERENCES
% ===========================================================================
{
  \small
  \bibliographystyle{ieee_fullname}
  \bibliography{references}
}

\end{document}
